\chapter{Fundamentos de Generación Text-to-Query y Consultas Híbridas}\label{chap:background}


Este capítulo presenta los conceptos necesarios para comprender la propuesta del presente trabajo, la cual será desarrollada en el Capítulo~\ref{chap:proposal}. La exposición avanza de manera progresiva, desde los fundamentos de los modelos de lenguaje hasta los mecanismos que permiten controlar y verificar las consultas generadas.

En primer lugar, se describen los \acp{LLM}, su forma de generar texto y los tipos de desafíos que pueden aparecer en este proceso (Sección~\ref{sec:modelos-lenguaje}). Luego, se revisan los modelos de datos relacional y de grafos de propiedad, así como su integración mediante el estándar \ac{SQL}/\ac{PGQ} (Sección~\ref{sec:modelos-datos}) y se introduce el problema de \emph{semantic parsing} y, dentro de este, la tarea de \emph{text-to-query} (Sección~\ref{sec:semantic-parsing}).

Posteriormente, se explica el papel de una representación intermedia tipada como puente entre la interpretación realizada por el modelo y la generación final de una consulta ejecutable (Sección~\ref{sec:ri}). Finalmente, se revisan las formas de verificación estática y dinámica que permiten evaluar la corrección de la consulta generada (Sección~\ref{sec:verificacion}).

Estos fundamentos sirven de base para el análisis del estado del arte y para la propuesta de esta investigación.

\section{Modelos de lenguaje y generación estructurada}
\label{sec:modelos-lenguaje}

Los \acp{LLM}, o modelos extensos de lenguaje, son modelos basados en arquitecturas \emph{transformer}~\cite{Vaswani2017Attention} y entrenados sobre grandes volúmenes de texto, los cuales en los últimos años se han convertido en una base importante para la investigación en el área de \emph{text-to-query} debido a su capacidad para interpretar preguntas en lenguaje natural y poder generar respuestas en formatos estructurados. Esto en adición a su flexibilidad para adaptarse a nuevos esquemas de bases de datos mediante aprendizaje en contexto \cite{Brown2020GPT3} les permite generar consultas sin requerir un reentrenamiento completo del modelo para cada dominio o base de datos. 

Esta capacidad de generalización se complementa con la versatilidad de los modelos preentrenados, los cuales pueden especializarse en la generación de consultas estructuradas mediante \emph{fine-tuning}, diseño avanzado de \emph{prompts} o enfoques híbridos que combinan ambas estrategias.

No obstante, esta capacidad no resuelve por sí sola el problema presentado, ya que estos modelos se caracterizan por generar texto de forma autorregresiva, es decir que se produce una secuencia de \emph{tokens} prediciendo cada elemento a partir de los anteriores, siendo este un proceso estadístico, que no incorpora de manera nativa las reglas gramaticales ni el esquema sobre el cual operar, lo que abre la puerta a la posibilidad de producir consultas que parezcan correctas desde el punto de vista del lenguaje natural, pero que contienen errores estructurales, referencias inexistentes o interpretaciones equivocadas de esquema. Fenómeno que en la literatura se conoce como \emph{alucinaciones}~\cite{Ji2023Hallucination}. 

\subsection{Arquitectura transformer y decodificación autorregresiva}
La arquitectura \emph{transformer}~\cite{Vaswani2017Attention} permite procesar una secuencia de entrada $x = (x_1, \ldots, x_n)$ y representar cada \emph{token} mediante vectores contextualizados. Para ello utiliza mecanismos de atención, los cuales permiten que el modelo relacione distintos elementos de la secuencia, incluso cuando se encuentran alejados entre sí. A diferencia de arquitecturas recurrentes, este procesamiento no depende de recorrer la secuencia paso a paso, sino de comparar los \emph{tokens} entre sí mediante capas de atención.

En tareas de \emph{text-to-query}, los \acp{LLM} suelen emplearse como modelos generativos autorregresivos. Esto significa que generan la salida $y = (y_1, \ldots, y_m)$ un \emph{token} a la vez, donde cada nuevo \emph{token} depende de los generados previamente y de la entrada original. Formalmente, esta generación puede expresarse como:
\[P(y \mid x) \;=\; \prod_{t=1}^{m} P(y_t \mid y_{<t},\, x),\]
donde $P(y_t \mid y_{<t}, x)$ representa la probabilidad del siguiente \emph{token} dado el prefijo ya generado $y_{<t}$ y la entrada $x$. En la práctica, esta probabilidad se obtiene mediante una capa \emph{softmax} sobre el vocabulario del modelo. A partir de esa distribución, el siguiente \emph{token} puede seleccionarse mediante estrategias como la maximización local o el muestreo.

Esta forma de generación tiene dos consecuencias importantes para \emph{text-to-query}. En primer lugar, el modelo decide cada paso en función de la probabilidad del siguiente \emph{token}, pero no garantiza por sí mismo que la consulta completa sea válida. Por ejemplo, una elección localmente probable puede introducir el nombre de una columna inexistente o una relación incorrecta entre tablas. En segundo lugar, el espacio de salida del modelo está definido por \emph{tokens}, no por las reglas formales del lenguaje de consulta. Por ello, aunque el modelo pueda producir una secuencia que se parezca a una consulta válida, no existe una garantía nativa de que dicha secuencia respete la gramática del lenguaje ni las restricciones del esquema objetivo.


\subsection{Prompting y aprendizaje en contexto}

Sobre esta base de generación, la aplicación a tareas específicas admite dos regímenes complementarios. El \emph{fine-tuning} supervisado modifica los parámetros del modelo mediante pares (pregunta, consulta) con el fin de sesgar $P(y \mid x)$ hacia el lenguaje objetivo. El \emph{aprendizaje en contexto}~\cite{Brown2020GPT3}, por su parte, mantiene los parámetros congelados y condiciona la generación en un \emph{prompt} que incluye instrucciones y, opcionalmente, ejemplos de la tarea; en este régimen, la distribución efectiva de generación depende de la composición del \emph{prompt}, no de una actualización del modelo.

Ambos regímenes comparten una limitación estructural clave: la salida es el resultado de un muestreo sobre una distribución que desconoce las restricciones formales del lenguaje de consulta y del esquema. Por ello, el estudio de la fidelidad estructural de la salida requiere mecanismos externos al propio \ac{LLM}.

\section{Modelos de datos y lenguajes de consulta}
\label{sec:modelos-datos}

Las consultas estudiadas en esta investigación se formulan sobre dos modelos de datos: el modelo relacional y el modelo de grafos de propiedad. Aunque ambos pueden representar información relacionada, lo hacen mediante estructuras distintas y con lenguajes de consulta diferentes. Por ello, esta sección presenta primero cada modelo por separado y luego explica su integración mediante el estándar \ac{SQL}/\ac{PGQ}.


\subsection{Modelo relacional y álgebra relacional}

El modelo relacional~\cite{Codd1970Relational} representa los datos mediante conjuntos de tuplas de aridad fija, donde cada posición se asocia a un atributo con dominio de valores. A su vez un \emph{esquema relacional} $\mathcal{S}_{\mathrm{rel}}$ agrupa estas relaciones y sus restricciones de integridad, como claves primarias, claves foráneas, restricciones de unicidad y dominios válidos.

El álgebra relacional proporciona los operadores formales para consultar este modelo. Entre sus operadores básicos se encuentran la selección $\sigma$, la proyección $\pi$, la unión natural $\bowtie$, la unión $\cup$, la diferencia $-$ y el renombramiento $\rho$. Estos operadores son cerrados sobre relaciones, por lo que sus resultados pueden combinarse nuevamente con otros operadores.

Esta base formal es relevante para esta investigación porque permite razonar sobre una consulta antes de ejecutarla. Por ejemplo, es posible verificar si una columna referenciada existe en el esquema, si una operación respeta las relaciones definidas entre tablas o si una composición de operadores es compatible con las restricciones disponibles. De este modo, el modelo relacional no solo define cómo se estructuran los datos, sino también un marco para validar consultas de forma estática.

\subsection{SQL como lenguaje declarativo}

SQL (\emph{Structured Query Language}), estandarizado en la serie ISO/IEC 9075, es el lenguaje declarativo dominante para consultar bases de datos relacionales. Una consulta en forma \texttt{SELECT-FROM-WHERE-GROUP BY-HAVING-ORDER BY} se corresponde, con una expresión del álgebra relacional extendida con agregaciones y ordenamiento. La distinción entre la consulta declarativa y el plan físico ejecutado por el motor es relevante para esta tesis: el objetivo de la generación es producir una expresión declarativa bien formada y aterrizada al esquema, en cambio la optimización y ejecución pasan a ser responsabilidad del \ac{RDBMS}.

\subsection{Grafos de propiedad y pattern matching}

Un \emph{grafo de propiedad}~\cite{Angles2017Graph} se define como una estructura
\[
G \;=\; (V,\, E,\, \lambda,\, \mathrm{src},\, \mathrm{tgt},\, \mathrm{prop}),
\]
donde $V$ es un conjunto finito de vértices, $E$ un conjunto finito de aristas, $\lambda: V \cup E \to 2^{\Sigma_L}$ asigna a cada elemento un conjunto de etiquetas de un alfabeto, $\Sigma_L$, $\mathrm{src}, \mathrm{tgt}: E \to V$ determinan los extremos de cada arista dirigida, y $\mathrm{prop}: (V \cup E) \times K \rightharpoonup D$ asocia a cada elemento y cada clave $k \in K$ un valor parcial en un dominio $D$.

La forma natural de consultar esta estructura es mediante \emph{pattern matching}. En lugar de describir una consulta como una combinación de tablas, se define un patrón $P$ que debe encontrarse dentro del grafo. Dicho patrón contiene variables, etiquetas y restricciones sobre propiedades. El resultado de la consulta corresponde a las coincidencias entre el patrón y el grafo, es decir, a las asignaciones de sus variables a vértices y aristas concretas que respetan las etiquetas y condiciones declaradas.

Dentro de este tipo de consultas, una clase importante son las \emph{regular path queries} (\acp{RPQ}), que permiten describir caminos mediante expresiones regulares sobre las etiquetas de las aristas. Estas consultas generalizan las búsquedas de vecindad, ya que no solo permiten consultar relaciones directas, sino también trayectorias de longitud variable que siguen una secuencia determinada de etiquetas~\cite{Angles2017Graph}.

Esta formalización es relevante para la verificación de consultas sobre grafos ya que permiten formular reglas de tipado análogas a las del modelo relacional: toda referencia en una consulta sobre grafos debe corresponder a una etiqueta, un tipo o una propiedad declarada en la especificación de $G$.

\subsection{El estándar SQL/PGQ como integración}

El estándar ISO/IEC 9075-16:2023~\cite{ISO9075_16_2023} introduce \ac{SQL}/\ac{PGQ} como una extensión de SQL orientada a la consulta de grafos de propiedad definidos sobre tablas relacionales existentes. Como su construcción principal tenemos la expresión:
\[ \texttt{GRAPH\_TABLE(graph MATCH pattern COLUMNS (...))} \]
La cual evalua un patrón de grafo y devuelve una relación derivada que puede usarse en conjunto con las cláusulas tradicionales de SQL. 

Esta integración resulta relevante dado que une dos formas distintas de representar y consultar datos mantieniendo, por un lado, el modelo relacional (basado en tablas, atributos y restricciones); e integrandolo con el modelo de grafos de propiedad (expresado mediante patrones de vértices, aristas, etiquetas y propiedades).

Esta combinación introduce un reto adicional en el campo investigado ya que no solo se debe generar consultas SQL válidas, sino también patrones de grafo coherentes, así como se debe preservar la consistencia entre ambos niveles. Es por esto que una consulta SQL/PGQ no puede validarse únicamente revisando cada componente por separado, sino que también es necesario verificar que la interacción entre ambos sea estructural y semánticamente correcta.

\section{Semantic parsing y generación \emph{text-to-query}}
\label{sec:semantic-parsing}

\subsection{\emph{Semantic parsing}: definición y marco general}

El \emph{semantic parsing}~\cite{ZelleMooney1996} es la tarea de transformar un enunciado en lenguaje natural $u$ en una representación formal $z$ que pueda ser ejecutada sobre un entorno $\mathcal{E}$. Este entorno puede ser, por ejemplo, una base de datos, un grafo de conocimiento o incluso un sistema físico. El objetivo es que la ejecución de dicha representación formal, denotada como $\llbracket z \rrbracket_{\mathcal{E}}$, produzca la respuesta esperada para la pregunta original.

La representación formal $z$ puede adoptar distintas formas, dependiendo del dominio y del tipo de sistema. Por ejemplo, puede tratarse de expresiones en cálculo, lenguajes específicos de dominio, consultas SQL, consultas sobre grafos como \emph{Cypher}, entre otros. Sea cual sea el caso, este problema requiere relacionar tres componentes: el enunciado en lenguaje natural, el entorno sobre el cual se ejecutará la consulta y el mecanismo que traduce el primero en una representación formal válida.

\subsection{\emph{Text-to-SQL} como caso particular y extensión a \emph{text-to-query}}

La tarea \emph{text-to-SQL} es un caso particular de \emph{semantic parsing}. En este caso, la representación formal pertenece al lenguaje SQL y el entorno de ejecución es una base de datos relacional. La entrada del sistema puede expresarse como el par $(u, \mathcal{S}_{\mathrm{rel}})$, donde $u$ es la pregunta en lenguaje natural y $\mathcal{S}_{\mathrm{rel}}$ es el esquema de la base de datos. La salida esperada es una consulta $q$ que, al ejecutarse sobre la base de datos, produzca la respuesta correcta.

La presente tesis aborda una extensión de este problema hacia \emph{text-to-query} sobre \ac{SQL}/\ac{PGQ}. En este escenario, la consulta no se limita únicamente a operaciones relacionales, sino que también puede incluir patrones de grafo y combinaciones entre ambos modelos. Por ello, la entrada se define como un par $(u, \mathcal{S})$, donde $\mathcal{S}$ representa un esquema dual:

\[
\mathcal{S} \;=\; (\mathcal{S}_{\mathrm{rel}},\, \mathcal{S}_{\mathrm{grafo}}).
\]

Aquí, $\mathcal{S}_{\mathrm{rel}}$ corresponde al componente relacional, mientras que $\mathcal{S}_{\mathrm{grafo}}$ describe el componente de grafo de propiedad. La salida esperada es una consulta capaz de combinar operadores SQL tradicionales con patrones de grafo definidos bajo \ac{SQL}/\ac{PGQ}.

\section{Representaciones intermedias}
\label{sec:ri}

\subsection{Noción formal de representación intermedia}

Una \emph{representación intermedia} (\ac{RI}) es un lenguaje formal $\mathcal{L}_{\mathrm{RI}}$ dotado de una gramática $\mathcal{G}_{\mathrm{RI}}$, un sistema de tipos $\mathcal{T}_{\mathrm{RI}}$ y una función de compilación
\[
\kappa: \mathcal{L}_{\mathrm{RI}} \to \mathcal{L}_{\mathrm{objetivo}}
\]
hacia el lenguaje de consulta efectivo. Su rol en un pipeline de generación es mediar entre la salida del componente estadístico y la consulta ejecutable: el modelo produce, o es guiado a producir, instancias bien tipadas de $\mathcal{L}_{\mathrm{RI}}$; el compilador $\kappa$ genera determinísticamente la consulta final.

\subsection{Propiedades deseables}

Una \ac{RI} adecuada para \emph{text-to-query} satisface, idealmente, las siguientes propiedades, identificadas a partir de los diseños de \acp{RI} previas para parsing semántico \cite{guo2019irnet,gan2021natsql,herzig2021unlocking,nie2022graphqir} y de la formalización de los lenguajes objetivo \cite{francis2023gpc,deutsch2022graph}:

\begin{description}
\item[Abstracción sintáctica.] $\mathcal{L}_{\mathrm{RI}}$ suprime decisiones sintácticas del lenguaje objetivo que no tienen contraparte directa en la formulación en lenguaje natural, tales como alias, el orden de cláusulas, la elección entre \texttt{JOIN} y subconsulta, o la presencia explícita de \texttt{FROM}, \texttt{GROUP BY} y \texttt{HAVING}. La motivación sigue la línea de SemQL \cite{guo2019irnet} y, más radicalmente, NatSQL \cite{gan2021natsql}, que prescinden de operadores y cláusulas SQL ``raramente encontradas en las descripciones del usuario''. Como consecuencia, consultas objetivo semánticamente equivalentes tienden a colapsar en una representación canónica en $\mathcal{L}_{\mathrm{RI}}$, reduciendo el espacio de salida que el modelo debe predecir.

\item[Tipado composicional.] La buena formación de una instancia se decide mediante un sistema de tipos estático: cada constructor declara los tipos de sus operandos y el tipo del término resultante se deriva de los de sus subtérminos. Esta propiedad sigue la tradición de los cálculos formales para lenguajes de consulta de grafos, en particular GPC \cite{francis2023gpc}, donde un conjunto de reglas de tipado garantiza que toda expresión bien tipada admite un esquema $\mathrm{sch}(\xi)$ derivable composicionalmente. El tipado opera puramente sobre la estructura de la \ac{RI} y el esquema de la base de datos, sin requerir evaluación.

\item[Compilación dirigida por sintaxis.] La función de compilación $\kappa : \mathcal{L}_{\mathrm{RI}} \times \mathcal{S} \rightarrow \mathcal{L}_{\mathrm{obj}}$, donde $\mathcal{S}$ es el esquema, es total sobre instancias bien tipadas y se define por inducción estructural sobre la \ac{RI}. La indeterminación propia de la generación neuronal queda así confinada al espacio $\mathcal{L}_{\mathrm{RI}}$, en línea con el principio de \emph{deterministic inference} de IRNet \cite{guo2019irnet} y con la separación entre parser semántico y compilador \emph{source-to-source} adoptada por GraphQ IR \cite{nie2022graphqir}. La dependencia explícita del esquema es necesaria porque parte de la información elidida en la \ac{RI} (e.g.\ inferencia de \texttt{GROUP BY}, resolución de claves foráneas para \texttt{JOIN}) se reconstruye consultando $\mathcal{S}$. % LLM-like: cliché-LLM:"en línea con"

\item[Expresividad adecuada.] $\mathcal{L}_{\mathrm{RI}}$ debe cubrir las construcciones del lenguaje objetivo que aparecen en la distribución de tareas considerada. Para SQL/PGQ \cite{deutsch2022graph}, esto exige soporte simultáneo para las tres capas conceptuales del estándar \cite{bonifati2025expressiveness}: (i) emparejamiento de patrones de grafo, (ii) operadores del álgebra relacional, y (iii) composición híbrida mediante \texttt{GRAPH\_TABLE}. Las \acp{RI} existentes cubren parcialmente alguno de estos fragmentos, como podrían ser SemQL/NatSQL para SQL relacional \cite{guo2019irnet,gan2021natsql}y GraphQ IR para Cypher/SPARQL/KoPL \cite{nie2022graphqir}; pero, hasta donde pudo verificarse, no existe una \ac{RI} que cubra el fragmento híbrido relacional-grafo de SQL/PGQ.
\end{description}

\section{Verificación de consultas generadas}
\label{sec:verificacion}

Si bien el objetivo es la producción de una cadena de consulta, el proceso de \emph{text-to-query} no termina ahí, ya que es necesario evaluar si la consulta generada es consistente con el esquema, si puede ejecutarse correctamente y si su comportamiento es coherente frente a la pregunta original, a este proceso le denominamos verificación y en esta sección detallaremos tres niveles y como se aplican al problema investigado.

\subsection{Verificación estática}

La verificación estática designa el análisis de una consulta, o de su \ac{RI}, sin ejecutarla, mediante la aplicación de reglas sobre su estructura y el esquema. En el contexto de \emph{text-to-query}, los chequeos estáticos relevantes incluyen: existencia de tablas, columnas, etiquetas y propiedades referenciadas; consistencia de tipos en operadores y condiciones; satisfacibilidad trivial de restricciones, como detección de condiciones contradictorias evidentes; y, en consultas híbridas, coherencia entre los componentes relacional y de grafo.

Formalmente, la verificación estática puede modelarse como un conjunto de predicados $\Phi = \{\varphi_1, \ldots, \varphi_k\}$ sobre consultas y esquemas. Una consulta $q$ se considera \emph{estáticamente válida} respecto a un esquema $\mathcal{S}$ si se cumple $\bigwedge_{i=1}^{k} \varphi_i(q,\, \mathcal{S})$. La decidibilidad local de cada $\varphi_i$ hace de la verificación estática un mecanismo de costo constante o lineal respecto al tamaño de $q$, y por ello resulta aplicable en un bucle de refinamiento sin comprometer la eficiencia global.

\subsection{Verificación dinámica}

La verificación dinámica ejecuta la consulta en un entorno controlado y analiza la respuesta del motor. El motor actúa como \emph{oráculo parcial}: distingue consultas sintáctica y referencialmente válidas de aquellas que disparan errores de ejecución, pero no decide, en general, la corrección semántica de una consulta ejecutable. Las señales útiles en este régimen incluyen los errores tipados del motor, la cardinalidad anómala del resultado, entendida como un resultado vacío cuando se espera contenido o viceversa, la presencia inesperada de valores nulos y las desviaciones respecto a respuestas de referencia cuando estas existen.

\subsection{Testing metamórfico}

El \emph{testing metamórfico}~\cite{ChenTY1998Metamorphic} extiende la verificación dinámica más allá de lo explícito. La idea central consiste en identificar relaciones entre entradas y salidas que deben preservarse bajo transformaciones conocidas, denominadas \emph{relaciones metamórficas}. En \emph{text-to-query}, relaciones como la invariancia del resultado bajo renombramiento de alias, la inclusión del resultado de una consulta con \texttt{LIMIT} en el de la misma consulta sin dicha cláusula, o la preservación de cardinalidades bajo reformulaciones semánticamente equivalentes, ofrecen un mecanismo para detectar inconsistencias sin conocer de antemano la respuesta correcta. La violación de una relación metamórfica evidencia un error; su satisfacción, sin embargo, no certifica corrección, sino que acota el espacio de hipótesis de error.
